\icode[magenta]{E-MD-CA-L4-261113-002}
某国际机场某货物传送装置实物简化图如图\(a\)，该装置由传送带\(ABCD\)及固定挡板\(CDEF\)组成，固定挡板与传送带上表面垂直，传送带上表面\(ABCD\)与水平台面的夹角\(\theta =\)\,\qty{37}{\degree}。传送带匀速转动时，工作人员将正方体货物从\(D\)点由静止释放，在\(L =\)\,\qty{10}{m}处取下货物，货物运动时的剖面图如图\(b\)所示，已知传送带匀速运行的速度\(v =\)\,\qty{1}{m/s}，货物质量\(m =\)\,\qty{10}{kg}，其底部与传送带上表面\(ABCD\)的动摩擦因数为\(\mu_1 = 0.5\)，其侧面与挡板\(CDEF\)的动摩擦因数为\(\mu_2 = 0.25\)。（重力加速度\(g =\)\,\qty{10}{m/s^2}	，不计空气阻力）
\begin{enumerate}
	\item 求传送带上表面对货物的摩擦力大小\(f_1\)和挡板对货物的摩擦力大小\(f_2\)；
	\item 货物在传送带上运动的时间\(t\);
	\item 货物若传送到\(s =\)\,\qty{0.2}{m}时，传送带由于故障突然停止，工作人员待货物安全停止时立即进行维修，用时\(T =\)\,\qty{30}{s}，传送带恢复正常，忽略传送带加速至恢复正常的时间，求由于故障，传送货物耽误的时间\(\Delta t\)。（保留两位小数）
\end{enumerate}

\begin{tikzpicture}[>=stealth,font=\scriptsize]
	\draw (0,-.4)--(5,-.4) (5.24,-.2)--(6.44,1.4)--(6.44,.2);
	\draw[fill=teal!60] (0,-.4) rectangle(5,-.8) (5,-.8)--(6.44,-.2)--(6.44,.2)--(5,-.4)--(5,-.8);
	\draw (6.2,1.6) arc(90:0:.24 and .2);
	\draw[fill=white] (0,-.2) ellipse(.24 and .2) (5,-.2) ellipse(.24 and .2);
	\draw[fill=teal!60] (0,0)node[below]{\(D\)}--(5,0)node[below]{\(C\)}--(6.2,1.6)node[above]{\(B\)}--(1.2,1.6)node[above]{\(A\)}--(0,0);
	\draw[fill=teal!40] (0,0)--(5,0)--(4.7,.4)node[above]{\(F\)}--(-.3,.4)node[above]{\(E\)}--(0,0);
	\node at(5.5,-.05){\(\theta\)};
	\draw[->] (2.6,1)--(3.6,1)node[right]{\(v\)};
	\node at(3,-1){图\(a\)};
\end{tikzpicture}
\hspace{2em}
\begin{tikzpicture}[>=stealth,font=\scriptsize]
	\begin{scope}[rotate=-37]
		\draw[fill=teal!60] (0,0) rectangle(-3,.2);
		\node at(-3,.4)[above]{\(A\)};
		\draw[very thick] (0,0)--(0,1.2)node[right]{\(E\)};
		\draw[fill=orange!60] (0,.2) rectangle(-.4,.6);
	\end{scope}
	\draw (-.3,0) arc(180:143:.3);
	\draw[thick,dashed,cyan] (-2.5,0)--(.8,0);
	\node at(-.6,.2){\(\theta\)};
	\node at(.3,.2){\(D\)};
	\node at(-.8,-.5){图\(b\)};
\end{tikzpicture}
\begin{answer}
    （1）\qty{40}{N}\quad \qty{15}{N}（2）\qty{10.2}{s}（3）\qty{30.29}{s}
\end{answer}